% !TeX program = xelatex
% !TeX encoding = utf8
% !TeX spellcheck = uk_UA
% !BIB program = biber
\documentclass{Syllabus}
\addbibresource{d:/Projects/LaTeX/MyPackage/MyBase.bib}

% ================================================
%                Дані  дисципліни
% ================================================
\logokaf{logokaf.png}
\creator{доцент, к.ф.-м.н., доцент Пономаренко Сергій Миколайович}
\lecturer[s.ponomarenko@kpi.ua]{доцент, к.ф.-м.н., доцент Пономаренко Сергій Миколайович}
\practicer[]{доцент, к.ф.-м.н., доцент Пономаренко Сергій Миколайович}
\disciplinename{Фізика плазми} % Назва дисципліни
\disciplinecode{Ф-6}
\kursurl{https://classroom.google.com/u/2/c/NjU4MTMxNzc0MTkw}
\shedule{http://ipt.kpi.ua/navchalnij-protses}
% -------------- Номера протоколів ---------------
\methodcomday{29} % Число ухвалення метод момісією
\methodcommonth{червня} % місяць ухвалення метод момісією
\methodcomyear{2023}
\methodcomnum{6} % Номер протокола ухвалення метод момісією
\kafday{15} % Число ухвалення кафедрою
\kafmonth{червня} % місяць ухвалення кафедрою
\kafyear{2023}
\kafnum{{6}}% Номер протокола ухвалення кафедрою
% ---------------- Syllabus head -----------------
\speciality{105 Прикладна фізика та наноматеріали}
\kafname{прикладної фізики}
% ---------------- Syllabus head -----------------
\efield{10 Природничі науки}
\teachingprogramm{Прикладна фізика}
%\cycle{Цикл загальної підготовки}
\level{Другий (магістерський)} % освітньо-кваліфікаційний рівень
\status{Вибіркова}
%-------- Семестр 3 --------------------
\course{1}
\semestr{весняний}
\controltype{екзамен}
\AudDistribHours%
{5}%
{36}%
{36}%
{0}%


\MKR{1}
\DKR{}
\RGR{}
\Referat{1}
% ================================================


\def\zsrs{\textit{Завдання на СРС:}}
\def\lit{\textit{Опрацювати:\ }}
\def\probl{\textit{Задачі:\ }}
\def\lab{\textit{Лабораторний практикум:\ }}


\begin{document}
\printhead


\syllabuschapter{Реквізити навчальної дисципліни}
\bluetableprint

\syllabuschapter{Програма навчальної дисципліни}

\section{Опис навчальної дисципліни, її мета, предмет вивчання та результати навчання}


Метою курсу є підвищення загального рівня знань студентів, введення їх у коло явищ, які розглядає фізика плазмі, ознайомлення їх з основними термінами, поняттями, методами і ідеями, що були розвинуті в цій науці, підготовка їх до самостійної роботи із оригінальною науковою літературою.

Предметом вивчення є характерні явища в системах заряджених частинок: особливості колективного руху заряджених частинок завдяки їх електромагнітній взаємодії, та особливості індивідуального руху окремих заряджених частинок в електричних та магнітних полях.


Після засвоєння навчальної дисципліни студенти мають продемонструвати такі \emph{результати} навчання:

\begin{enumerate}
	\item[\bfseries знання:] загальних і характерних властивостей плазми як системи заряджених частинок з далекодійним характером взаємодії; основні поняття про природу елементарних та колективних процесів в плазмі, явищ переносу в плазмі за відсутності та наявності магнітного поля; особливості руху заряджених частинок в магнітному полі;  фізичні механізми електричного пробою та основні властивості плазми газового електричного розряду
	\item[\bfseries уміння:] орієнтуватися в термінології, основах математичного формалізму модельних концепцій; виявляти і прогнозувати специфічну поведінку систем заряджених частинок, що
	      зумовлюється далекодійною кулонівською та електромагнітною взаємодією.
	\item[\bfseries досвід:] вільно орієнтуватися на якісному й кількісному рівні в основних фізичних явищах, пов’язаних з проявами електромагнітного поля в плазмі, виробити навички практичного використання засвоєних знань, методів і підходів у подальшому навчанні та професійній діяльності.
\end{enumerate}

Згідно з вимогами освітньо-професійної програми студенти після засвоєння навчальної дисципліни <<\discipline>> мають продемонструвати такі результати навчання:

\subsection*{Загальні компетентності ОПП}

\begin{enumerate}
	\item [ЗК 1:] Здатність застосовувати знання у практичних ситуаціях.
	\item [ЗК 2:] Знання та розуміння предметної області та розуміння професійної діяльності.
	      %\item [ЗК 3:] Здатність спілкуватися державною мовою як усно, так і письмово.
	      %\item [ЗК 4:] Здатність спілкуватися іноземною мовою.
	      %\item [ЗК 5:] Здатність використання інформаційних і комунікаційних технологій.
	\item [ЗК 6:] Здатність проведення досліджень на відповідному рівні.
	      %\item [ЗК 7:] Здатність до пошуку, оброблення та аналізу інформації з різних джерел.
	      %\item [ЗК 8:] Здатність застосовувати навички міжособистісної взаємодії.
	      %\item [ЗК 9:] Здатність працювати автономно.
	      %\item [ЗК 10:] Здатність здійснювати безпечну діяльність.
	      %\item [ЗК 11:] Здатність реалізувати свої права і обов’язки як члена суспільства, усвідомлювати цінності громадянського (вільного демократичного) суспільства та необхідність його сталого розвитку, верховенства права, прав і свобод людини і громадянина в Україні.
	      %\item [ЗК 12:] Здатність зберігати та примножувати моральні, культурні, наукові цінності і досягнення суспільства на основі розуміння історії та закономірностей розвитку предметної області, її місця у загальній системі знань про природу і суспільство та  розвитку суспільства, техніки і технологій, використовувати різні види та форми рухової активності для активного відпочинку та ведення здорового способу життя.
	      %\item [ЗК 13:] Здатність критично оцінювати результати своєї діяльності в професійній сфері, навчанні і приймати обґрунтовані рішення з урахуванням наукових, соціальних, етичних, правових, економічних аспектів.
	      %\item [ЗК 14:] Здатність продовжувати навчання із значним ступенем самостійності
\end{enumerate}

\subsection*{Фахові компетентності ОПП}

\begin{enumerate}
	%\item [ФК 1:] Здатність брати участь у плануванні та виконанні наукових та науково-технічних проектів.
	\item [ФК 2:] Здатність брати участь у плануванні і виконанні експериментів та лабораторних досліджень властивостей фізичних систем, фізичних явищ і процесів, обробленні й презентації їхніх результатів.
	\item [ФК 3:] Здатність брати участь у виготовленні експериментальних зразків, інших об'єктів дослідження.
	      %\item [ФК 4:] Здатність брати участь у впровадженні результатів досліджень та розробок.
	      %\item [ФК 5:] Здатність до постійного розвитку компетентностей у сфері прикладної фізики, інженерії та комп’ютерних технологій.
	      %\item [ФК 6:] Здатність використовувати сучасні теоретичні уявлення в галузі фізики для аналізу фізичних систем.
	      %\item [ФК 7:] Здатність використовувати методи і засоби теоретичного дослідження та математичного моделювання в професійній діяльності.
	      %\item [ФК 8:] Здатність працювати в колективах виконавців, у тому числі в міждисциплінарних проектах.
	      %\item [ФК 9:] Здатність брати участь у роботах з проведення наукових досліджень властивостей явищ і процесів у фізичній та/або біофізичній, енергетичній системі, та зі складання наукових звітів з виконаних робіт.
	      %\item [ФК 10:] Здатність розуміти та застосовувати апарат спеціальних розділів математики для розв’язання проблем прикладної фізики, моделювати фізичні процеси і системи, використовуючи статистичні та стохастичні методи, комп’ютерну графіку, та представляти результати моделювання.
	      %\item [ФК 11:] Здатність використовувати знання основ професійно-орієнтованих дисциплін для виконання наукових досліджень, розв’язання практичних проблем прикладної фізики та для самостійного  опанування нових технологій, в тому числі із суміжних галузей, застосовувати отримані знання і практичні навички для прийняття інноваційних рішень при розв’язанні складних практичних задач або в навчанні, зокрема,  високих фізичних технологій та/або фізики живих систем та/або фізики енергетичних систем.
\end{enumerate}

\subsection*{Програмні результати навчання з ОПП}

\begin{enumerate}
	\item [ПРН 1:] Знати і розуміти сучасну фізику на рівні, достатньому для розв’язання складних спеціалізованих задач і практичних проблем прикладної фізики.
	      %\item [ПРН 2:] Застосовувати сучасні математичні методи для побудови й аналізу математичних моделей фізичних процесів.
	\item [ПРН 3:] Застосовувати ефективні технології, інструменти та методи експериментального дослідження властивостей речовин і матеріалів, включаючи наноматеріали, при розв’язанні практичних проблем прикладної фізики.
	\item [ПРН 4:] Застосовувати фізичні, математичні та комп'ютерні моделі для дослідження фізичних явищ, розробки приладів і наукоємних технологій.
	      %\item [ПРН 5:] Вибирати ефективні методи та інструментальні засоби проведення досліджень у галузі прикладної фізики.
	      %\item [ПРН 6:] Відшуковувати необхідну науково-технічну інформацію в науковій літературі, електронних базах, інших джерелах, оцінювати надійність та релевантність інформації.
	      %\item [ПРН 7:] Класифікувати, аналізувати та інтерпретувати науково-технічну інформацію в галузі прикладної фізики.
	      %\item [ПРН 8:] Вільно спілкуватися з професійних питань державною та англійською мовами усно та письмово.
	      %\item [ПРН 9:] Презентувати результати досліджень і розробок фахівцям і нефахівцям, аргументувати власну позицію.
	      %\item [ПРН 10:] Планувати й організовувати результативну професійну діяльність індивідуально і як член команди при розробці та реалізації наукових і прикладних проєктів.
	      %\item [ПРН 11:] Знати цілі сталого розвитку та можливості своєї професійної сфери для їх досягнення, в тому числі в Україні.
	      %\item [ПРН 12:] Розуміти закономірності розвитку прикладної фізики, її місце в розвитку техніки, технологій і суспільства, у тому числі в розв'язанні екологічних проблем.
	      %\item [ПРН 13:] Оцінювати фінансові, матеріальні та інші витрати, пов’язані з реалізацією проектів у сфері прикладної фізики, соціальні, екологічні та інші потенційні наслідки реалізації проектів.
	\item [ПРН 14:] Обирати та використовувати методи та засоби дослідження структури, складу та речовин і матеріалів.
	      %\item [ПРН 15:] Знання основ методології наукових досліджень в прикладній фізиці, технологіії  оформлення, презентації та захисту результатів наукових досліджень, вміння складати звіти з виконаних робіт.
	      %\item [ПРН 16:] Знання методів аналізу випадкових процесів, теорії ймовірності і математичної статистики, програмування, комп’ютерної графіки, прикладних програм і методів обчислень, методів розв’язання рівнянь математичної фізики, теорії функції комплексної змінної, тензорного аналізу, для розуміння сучасних фізичних теорій і розв’язання проблем прикладної фізики та моделювання процесів, що відбуваються в фізико-технічних системах.
	      %\item [ПРН 17:] Знання основ професійно-орієнтованих дисциплін спеціальності, зокрема  хімії, ядерної фізики, статистичної радіофізики та оптики, електродинаміки суцільних середовищ для розв’язання практичних проблем прикладної фізики, в т.ч. високих фізичних технологій та/або фізики живих систем та/або фізики енергетичних систем.
\end{enumerate}


\section{Пререквізити та постреквізити дисципліни (місце в структурно-логічній схемі\\ навчання за відповідною освітньою програмою)}

%Зазначається перелік дисциплін, або знань та умінь, володіння якими необхідні студенту (вимоги до рівня підготовки) для успішного засвоєння дисципліни (наприклад, «базовий рівень володіння англійською мовою не нижче А2»). Вказується перелік дисциплін які базуються на результатах навчання з даної дисципліни.

Для засвоєння матеріалу курсу <<\discipline>> студенти повинні знати курс фізики в рамках шкільної програми та засвоїти термінологію та поняття курсів:
\begin{enumerate}
	\item Математичний аналіз;
	\item Тензорний аналіз;
	\item Механіка;
	\item Термодинаміка та молекулярна фізика;
	\item Електрика та магнетизм;
	\item Оптика;
	\item Теорія поля;
	\item Електродинаміка суцільних середовищ.
\end{enumerate}

Також повинні вміти використовувати математичний апарат: операції з матрицями,  диференціювати, інтегрувати.

Засвоєні теоретичні знання та отримані практичні навички під час вивчення навчальної дисципліни «Фізика плазми» можна використовувати в навчальних дисциплінах, пов’язаних з вивченням явищ, технологій та методик, в яких використовуються або розглядаються явища пов’язані із плазмою, або з системами заряджених частинок (плазмові та іонно/електронно променеві технології та експериментальні методики/обладнання).


\section{Зміст навчальної дисципліни}

%Надається перелік розділів і тем всієї дисципліни.

\begin{Rozdil}
	\item Основні поняття про плазму.
    \item Рух заряджених частинок у плазмі.
    \item Гідродинаміка плазми.
	\item Термодинаміка плазми.
    \item Елементарні процеси в плазмі.
	\item Утримання плазми.
	\item Хвилі в однорідній рівноважній плазмі.
    \item Явища переносу в плазмі.
    \item Випромінювання плазми.
    \item Термоядерні реакції та проблема керованого термоядерного синтезу.
%    \item Плазма в промисловості.
%    \item Астрофізика плазми.
\end{Rozdil}

%Зазначається: базова (підручники, навчальні посібники) та додаткова (монографії, статті, документи, електронні ресурси) література, яку потрібно прочитати або використовувати для опанування дисципліни.
%
%Можна надати рекомендації та роз’яснення:
%\begin{itemize}
%\item де можна знайти зазначені матеріали (бібліотека, методичний кабінет, інтернет тощо);
%\item що з цього є обов’язковим для прочитання, а що факультативним;
%\item як саме студент/аспірант має використовувати ці матеріали (читати повністю, ознайомитись тощо);
%\end{itemize}
% зв’язок цих ресурсів з конкретними темами дисципліни.
%Бажано зазначати не більше п’яти базових джерел, які є вільно доступними, та не більше 20
%додаткових.


\section{Навчальні матеріали та ресурси}\label{sec:refsec}

Нижче наводиться перелік навчальних матеріалів та ресурсів для засвоєння матеріалу, розглядуваного на лекційних заняттях та для додаткового вивчення.

\nocite{%
	AnisimovPlasma2018,
	SytenkoMalnev,
	ZagorodnijMomot,
    ChenPlasma2016,
}

\printbibliography[category=Main, heading=subbibliography, title={Підручники та посібники}]
\addtocategory{Main}{%
	AnisimovPlasma2018,
	SytenkoMalnev,
	}

\printbibliography[category=Additional, heading=subbibliography, title={Додаткова}]
\addtocategory{Additional}{%
ChenPlasma2016,
ZagorodnijMomot,
}

%\printbibliography[category=Problems, heading=subbibliography, title={Задачники}]
%\addtocategory{Problems}{%
%
%}

\newpage


\syllabuschapter{Навчальний контент}

\section{Методика опанування навчальної дисципліни (освітнього компонента)}

%Надається інформація (за розділами, темами) про всі навчальні заняття (лекції, практичні,
%семінарські, лабораторні) та надаються рекомендації щодо їх засвоєння (наприклад, у формі
%календарного плану чи деталізованого опису кожного заняття та запланованої роботи).

\subsection*{Лекційні та практичні заняття}

Для поглибленого вивчення питань доцільно використовувати не лише рекомендований тематичний лекційний матеріал, а і ті джерела, що вказані в розділі \ref{sec:refsec}.

\DefTblrTemplate{contfoot-text}{default}{}
\DefTblrTemplate{conthead-text}{default}{}
\DefTblrTemplate{caption}{default}{}
\DefTblrTemplate{conthead}{default}{}
\DefTblrTemplate{capcont}{default}{}
\begin{longtblr}[]{
	colspec = {Q[l, m, fg=main]X[j,m,fg=main]Q[c,m,fg=main]Q[c,m,fg=main]Q[c,m,fg=main]},
    rowhead=1,
    row{1} = {bg=rowfill},
    hlines = {1-Z}{infotablecolor},
    vlines = {1-Z}{infotablecolor},
	}
    % --------------------------------------------------------------------------
	№           &  \SetCell{c} Назва теми лекції та перелік основних питань
    & Лекції & Практичні   & {Самостійне \\ опрацювання}                                                                   \\
    % --------------------------------------------------------------------------
	\rownumber. & \textbf{Основні поняття про плазму}.
    Класифікація видів плазми. Дебаївський радіус. Ленгмюрівська частота. Теплові флуктуації в плазмі.
    & + & + & -- \\
    % --------------------------------------------------------------------------
   	\rownumber. & \textbf{Рух частинок в плазмі}.
    Дрейфове наближення. Рух в однорідному магнітному полі. Електричний дрейф. Дрейф під дією малої сили. Гравітаційний дрейф. Градієнтний дрейф. Відцентровий дрейф. Поляризаційний дрейф. Дрейф у неоднорідному електричному полі. Проблема втримання плазми в магнітному полі. Магнітні пастки.  Нестійкості в плазмі.
    & + & + & -- \\
    % --------------------------------------------------------------------------
  	\rownumber. & \textbf{Гідродинаміка плазми}.
    Основне рівняння магнітної гідродинаміки. Однорідинна та дворідинна моделі. Вмороженість магнітного поля в плазму. Дифузія магнітного поля в плазмі.
    & + & + & -- \\
    % --------------------------------------------------------------------------
   	\rownumber. & \textbf{Термодинаміка плазми}.
    Кінетична теорія, рівняння стану, енергія плазми.
    & + & + & -- \\
    % --------------------------------------------------------------------------
	\rownumber. & \textbf{Елементарні процеси в плазмі}.
    Іонізація електронним ударом. Потрійна рекомбінація. Ступінчаста іонізація. Фоторекомбінація. Фотоіонізація. Резонансна перезарядка. Корональна рівновага.
    & + & + & -- \\
    % --------------------------------------------------------------------------
   	\rownumber. & \textbf{Хвильові явища в плазмі}.
    Хвилі Альфвена, хвилі звуку, поперечні хвилі, нестійкості.
    & + & + & -- \\
    % --------------------------------------------------------------------------
    \rownumber. & \textbf{Явища переносу в плазмі}.
    Кінетичні коефіцієнти. Теплопровідність плазми. Рівняння теплопровідності. Термосила. Провідність плазми. Перенесення тепла струмом. В'язкість плазми. Виділення тепла. Амбіполярна дифузія. Кінетичні коефіцієнти в магнітному полі. Бомівська дифузія. Узагальнений закон Ома. Ефект Холла.
    & + & + & -- \\
    % --------------------------------------------------------------------------
  	\rownumber. & \textbf{Випромінювання плазми}.
    Суцільний та лінійчатий спектр. Магніто-гальмівне випромінювання.      Рекомбінаційне випромінювання. Лінійчатий спектр. Доплерівське та Штарківське уширення ліній.
    & + & + & -- \\
    % --------------------------------------------------------------------------
    \rownumber. & \textbf{Термоядерні реакції}.
    Фізика ядерних реакцій. Паливні цикли. Кулонівський бар'єр. Параметри термоядерних реакцій. Запалювання термоядерної реакції. Керований термоядерний синтез.
    & - & + & + \\
    % -----------------------------------------------------------------------
%    \rownumber. & \textbf{Плазмові процеси в промисловості}.
%    Застосування плазми в технологіях, плазмові реактори.
%    & -- & -- & + \\
    % --------------------------------------------------------------------------
%    \rownumber. &  \textbf{Астрофізика плазми}.
%    Зірки, галактики, космічні промені, плазмові потоки в космосі.
%    & -- & -- & +
    % --------------------------------------------------------------------------
\end{longtblr}




\section{Самостійна робота студента}

Самостійна робота студентів має на меті розвиток творчих здібностей та активізація їх розумової діяльності, формування потреби безперервного самостійного поповнення знань та розвиток морально-вольових зусиль. Завданням самостійної роботи студентів є навчити студентів самостійно працювати з літературою, творчо сприймати навчальний матеріал і осмислювати його та формування навичок до щоденної роботи з метою одержання та узагальнення знань, умінь і навичок.


На самостійну роботу відводяться наступні види завдань:
\begin{enumerate}[label=$\bullet$]
	\item обробка і осмислення інформації, отриманої безпосередньо на заняттях;
    \item самостійне вивчення окремих тем або питань;
	\item виконання підготовчої роботи до практичних занять та до написання МКР;
	\item написання та підготовка до захисту реферату;
	\item підготовка до складання семестрового контролю.
\end{enumerate}




\syllabuschapter{Політика та контроль}

\section{Політика навчальної дисципліни (освітнього компонента)}

\subsection*{Відвідування занять}
Відвідування лекцій, а також відсутність на них, не оцінюється. Однак, студентам бажано відвідувати заняття, оскільки на них викладається теоретичний матеріал та розвиваються навички, необхідні для успішного складання \control у.

\subsection*{Пропущені контрольні заходи}

Результат модульної контрольної роботи для студента, який не з’явився на контрольний захід, є нульовим. У такому разі, студент має можливість написати модульну контрольну роботу, але максимальний бал за неї буде дорівнювати 50~\% від загальної кількості балів. Повторне написання модульної контрольної роботи не допускається.

\subsection*{Календарний рубіжний контроль}

Проміжна атестація студентів (далі --- атестація) є календарним рубіжним контролем. Метою проведення атестації є підвищення якості навчання студентів та моніторинг виконання графіка освітнього процесу студентами\footnote{Рейтингові системи оцінювання результатів навчання: Рекомендації до розроблення і застосування. Київ: КПІ ім. Ігоря Сікорського, 2018. 20 с.}.

\begin{center}
	\begin{tblr}{
        colspec={|Q[l,m,fg=main]|X[c,m,fg=main]|X[c,m,fg=main]|},
        row{1}={rowfill},
        hlines = {1-Z}{infotablecolor},
        vlines = {1-Z}{infotablecolor},
        }
	    Термін атестації        & {Перша атестація 8-й тиждень} & {Друга атестація 14-й тиждень}  \\
	    Критерій: поточний контроль & $\ge 20$~ балів               & $\ge 30$~ балів                \\
	\end{tblr}
\end{center}

Ці бали студенти отримують за результатами виконання домашніх завдань, які задаються в class\-room.

\subsection*{Академічна доброчесність}

Політика та принципи академічної доброчесності визначені у розділі 3 Кодексу честі Національного технічного університету України «Київський політехнічний інститут імені Ігоря Сікорського». Детальніше: \url{https://kpi.ua/code}.

\subsection*{Норми етичної поведінки}

Норми етичної поведінки студентів і працівників визначені у розділі 2 Кодексу честі Національного технічного університету України «Київський політехнічний інститут імені Ігоря Сікорського». Детальніше: \url{https://kpi.ua/code}.

\subsection*{Процедура оскарження результатів контрольних заходів}

Студенти мають можливість підняти будь-яке питання, яке стосується процедури контрольних заходів та очікувати, що воно буде розглянуто згідно із наперед визначеними процедурами (згідно <<Положення про систему забезпечення якості вищої освіти у Національному технічному університеті України «Київський політехнічний інститут імені Ігоря Сікорського>>, <<Положення про організацію навчального процесу>>).

\section{Види контролю та рейтингова система оцінювання результатів навчання (РСО)}

%Вказуються всі види контролю та бали за кожен елемент контролю, наприклад:
%Поточний контроль: експрес-опитування, опитування за темою заняття, МКР, тест тощо
%Календарний контроль: провадиться двічі на семестр як моніторинг поточного стану
%виконання вимог силабусу.
%
%Семестровий контроль: екзамен / залік / захист курсового проекту (роботи)
%Умови допуску до семестрового контролю: мінімально позитивна оцінка за індивідуальне
%завдання / зарахування усіх лабораторних робіт / семестровий рейтинг більше ХХ балів.

Видами контролю успішності засвоєння матеріалу дисципліни є  модульна контрольна робота (МКР), реферат та семестровий контроль.

\subsection*{Практичні заняття}

\pgfmathsetmacro{\pracBal}{2.5}
\pgfmathsetmacro{\pracNum}{8}
Практичні заняття оцінюються на основі розв'язаних практичних задач, які задаються в classroom. Обчислюється середній бал за всі здані практичні задачі, який множиться на ваговий коефіцієнт $0.2$. Максимальний бал за кожне практичне заняття \pracBal{} (з урахуванням вагового коефіцієнту), максимальний бал
\xinttheexpr\pracBal*\pracNum\relax.
%\pgfmathparse{\pracBal*\pracNum}\pgfmathprintnumber[precision=0]\pgfmathresult.

Бал за задачі, які були здані пізніше вказаного в classroom терміну, знижується на 20\%.

Задачі бажано оформлювати у форматі \texttt{PDF}, який створюється за допомогою \LaTeX.



\subsection*{Модульна контрольна робота}

\pgfmathsetmacro{\mkrBal}{10}

Модульна контрольна робота проводиться після завершення першої частини курсу <<\discipline>> проводиться протягом 2-х академічних годин на практичних заняттях. Вона складається з $3$-х задач і передбачає письмовий розв’язок задачі, подібних до тих, що розглядались на практичних заняттях та під час виконання домашніх робіт.

Критерії оцінювання модульної контрольної роботи:
\begin{enumerate}[label=$\bullet$]
	\item максимальна кількість балів за кожну задачу --- повна правильна відповідь, $95$\% інформації, там де треба наведено рисунки, позначення, є письмовий коментар щодо базових понять та законів, які використовуються під час розв’язку задачі,
	\item $75$\%  --- розв’язок правильний, не всі умови попереднього пункту виконано,
	\item $60$\%  --- наведено основні базові поняття для розв’язку, розв’язок неправильний.
	\item списані відповіді, які студент не може пояснити, не зараховуються.
\end{enumerate}


\subsection*{Реферат}

Тема роботи обирається студентом за власним бажанням і узгоджується з викладачем або обирається із списку (див. розділ~\ref{sec:referat}), який надає викладач. Тема роботи не повинна дублювати лекційний матеріал. Найбільш вітається робота, яка містить критичний аналіз оригінальних наукових публікацій. Допускається детальний аналіз тем, викладених в посібниках, але не освітлених в лекційному курсі. Оцінка визначається комплексом: вибір тематики, оформлення, творчий підхід. Фінальна оцінка складається на основі захисту --- публічної співбесіди з викладачем у присутності студентів. В обговоренні реферату можуть брати участь усі студенти групи. Реферат оцінюється за наступними критеріями Якість викладення результатів (макс. $30$ балів) Якість оформлення роботи (макс. $20$ балів). Захист реферату (макс. $50$ балів). Ваговий коефіцієнт $0.3$.

\pgfmathsetmacro{\referatBal}{30}



\subsection*{Умови допуску до \control у}

В таблиці наведені умови допуску до семестрового контролю.

\begin{center}\setcounter{magicrownumbers}{0}
	\begin{tblr}{
        colspec={|Q[c,m,fg=main]|Q[l,m,fg=main]|Q[l,m,fg=main]|},
        row{1}={rowfill},
        hlines = {1-Z}{infotablecolor},
        vlines = {1-Z}{infotablecolor},
    }
	    №      & Обов’язкова умова допуску до \control у & Критерій  \\
	    \rownumber & Потчний рейтинговий бал             & $\ge 30$  \\
	    \rownumber & МКР                                 & виконана  \\
	    \rownumber & Реферат                             & захищений \\
	\end{tblr}%
\end{center}

\subsection*{Семестровий контроль (\control)}

\pgfmathsetmacro{\controlBal}{40}

\Control~--- письмова відповідь на \emph{\color{main} два} теоретичних питання білету.

Оцінка за \control\ складається з рейтингу за семестр та оцінки за письмову відповідь напитання білету (максимум $40$ баллів).

Оцінка формується на основі наступних критеріїв:
\begin{enumerate}[label=$\bullet$]
	\item від $35$ до $40$ --- повна правильна відповідь, $95$\% інформації, наведено рисунки, позначення, є письмовий коментар щодо базових понять та законів, формулювання та терміни точні, терміни роз’яснено, повна правильна відповідь на уточнюючі запитання
	\item від $25$ до $35$ --- правильна відповідь, 80\% інформації, наведено рисунки, позначення, є письмові коментарі щодо базових понять та законів, формулювання та терміни по суті правильні але не повні, терміни роз’яснено, правильна відповідь на уточнюючі запитання
	\item від $20$ до $25$ --- по суті правильна але неповна відповідь, $70$\% інформації, наведено рисунки та позначення, відсутні письмові коментарі щодо базових понять та законів, формулювання та терміни по суті правильні але не повні, терміни не роз’яснено, правильна відповідь на більшість уточнюючих запитання
	\item від $10$ до $20$ --- відповідь неповна, $50$\% інформації, не наведено рисунки та позначення, відсутні письмові коментарі щодо базових понять та законів, формулювання та терміни в основному правильні але не повні, терміни не роз’яснено, відповіді на уточнюючі запитання не повні
	\item від $0$ до $10$ --- відповідь неповна, $30$\% інформації, не наведено рисунки та позначення, відсутні письмові коментарі щодо базових понять та законів, формулювання та терміни в основному не повні, терміни не роз’яснено, відповіді на уточнюючі запитання не повні або відсутні
\end{enumerate}


Остаточна оцінка $\mathbf{R}$ є сумою рейтингових балів отриманих за поточний контроль та балів отриманих на \control і після співбесіди зі студентом.


\begin{center}\setcounter{magicrownumbers}{0}
	\begin{spreadtab}{{tblr}
    {colspec={Q[c,m,fg=main]Q[l,m,fg=main]Q[c,m,fg=main]Q[c,m,fg=main]Q[c,m,fg=main]},
    row{1}={rowfill},
    hlines = {1-Z}{infotablecolor},
    vlines = {1-Z}{infotablecolor},
    }
    }
	    @ №                       & @ Контрольний захід          & @ Бал       & @ Кількість                    & @ Всього                    \\
	    @ \rownumber              & @ Практичні заняття          & \pracBal    & \pracNum                       & \STcopy{v}{c2*d2}           \\
	    @ \rownumber              & @ Модульна контрольна робота & \mkrBal     & 1                              &                              \\
	    @ \rownumber              & @ Реферат                    & \referatBal & 1                              &                              \\
	    @ \rownumber              & @ \Control                   & \controlBal & 1                              &                              \\
	    @\SetCell[c=4]{h,c}Всього &                              &                                &             & sum(e1:[0,-1]) \\
	\end{spreadtab}
\end{center}






Таблиця відповідності рейтингових балів оцінкам за університетською шкалою.
\begin{center}
	\begin{tblr}{
        colspec={Q[c,m,fg=main]Q[c,m,fg=main, 0.25\linewidth]},
        row{1}={rowfill},
        hlines = {1-Z}{infotablecolor},
        vlines = {1-Z}{infotablecolor},
    }
	            Значення рейтингу              & Оцінка ECTS  \\
	       $95 \le \mathbf{R} \le 100$         &   відмінно   \\
	         $85 \le \mathbf{R} < 95$          &  дуже добре  \\
	         $75 \le \mathbf{R} < 85$          &    добре     \\
	         $65 \le \mathbf{R} < 75$          &  задовільно  \\
	         $60 \le \mathbf{R} < 65$          &  достатньо   \\
	            $ \mathbf{R} < 60$             & незадовільно \\
	    {Не здана МКР,\\ не захищений реферат} & не допущено  \\
	\end{tblr}%
\end{center}

%\clearpage

\section{Додаткова інформація з дисципліни}

\subsection*{Перелік рекомендованих тем рефератів}\label{sec:referat}
\begin{enumerate}
\item Газові розряди: їх види, фізика процесів та перспективи застосування.
\item Плазмові двигуни для космічних апаратів: принципи роботи, типи, перспективи.
\item Досягнення в діагностиці плазми: методи, прилади, принципи вимірювання.
\item Плазмові лазери: особливості генерації, посилення та керування плазмовим випромінюванням.
\item Плазмова обробка матеріалів: методи, застосування, ефекти.
\item Плазмова нанотехнологія: синтез, модифікація та застосування наночастинок і наноструктур з використанням плазми.
\item Плазми в електроніці: створення ефективних сенсорів та пристроїв зберігання інформації.
\item Використання плазми для вирішення проблем екології.
\item Плазмові джерела електричної енергії: принципи, проблеми, сучасний стан розвитку та перспективи.
\item Термоядерна енергетика: принципи, проблеми, перспективи.
\item Плазмова біологія та медицина: вплив плазми на живі організми, біомедичне застосування плазми.
\item Плазма в астрофізиці: сонячний вітер та його вплив на магнітосферу планет.
\end{enumerate}











%\begin{itemize}
%\item перелік питань, які виносяться на семестровий контроль (наприклад, як додаток до силабусу);
%\item  можливість зарахування сертифікатів проходження дистанційних чи онлайн курсів за відповідною тематикою;
%\item  інша інформація для студентів/аспірантів щодо особливостей опанування навчальної дисципліни.
%\end{itemize}

%\section{Додаткова інформація з дисципліни <<\discipline>>}


%•	перелік питань, які виносяться на семестровий контроль (наприклад, як додаток до силабусу);
%•	можливість зарахування сертифікатів проходження дистанційних чи онлайн курсів за відповідною тематикою;
%•	інша інформація для студентів/аспірантів щодо особливостей опанування навчальної дисципліни.

%\subsection*{Методи навчання}
%
%Методом проведення лекційних занять є наочний метод, який включає демонстрацію ключових дослідів та використання ілюстративного матеріалу; проблемний виклад, який дає уявлення про набуття нових знань на основі відомих фактів; бесіди і дискусії, для встановлення діалогу з аудиторією та залучення їх для активного осмислення інформації.
%
%Методами проведення практичних занять є метод тренінгу та репродуктивний метод. Для лабораторних занять використовуються частково-пошукові та дослідницькі методи.

%По всім видам занять проводяться дистанційні консультації.

%\subsection*{Перелік питань, які виносяться на усний \control}
%
%\loaditemsfromfile[enumerate]{ElMagQ.dat}
%
\phantom{text}
\vfill
\thispagestyle{empty}
\printrequisites
\end{document}



